\documentclass[10pt,twocolumn,letterpaper]{article}

\usepackage{cvpr}
\usepackage{times}
\usepackage{epsfig}
\usepackage{graphicx}
\usepackage{amsmath}
\usepackage{amssymb}
\usepackage{booktabs}
\usepackage{multirow}

\title{Advanced Video Object Removal: A Hybrid Framework of ProPainter and Selective Diffusion Inpainting}

\author{
Author Name\\
Institution\\
{\tt\small email@domain.com}
}

\begin{document}

\maketitle

\begin{abstract}
This report details an advanced pipeline for video object removal that integrates Stable Diffusion XL (SDXL) into a traditional patch-based propagation framework (ProPainter). While the baseline ProPainter effectively maintains temporal consistency, it often struggles to synthesize plausible, high-frequency textures for large, complex missing regions. To mitigate this, we propose sparsely applying SDXL specifically to these challenging sub-regions on selected keyframes. However, this hybrid architecture introduced two major artifacts: 1) spatial "halos" (white edges) caused by VAE boundary padding and morphological hard-cuts, and 2) temporal flickering resulting from independent frame-by-frame generative synthesis. We successfully resolve the spatial halos by implementing a Distance Transform Alpha blending mechanism, yielding seamless integration. To address temporal flickering, we extensively explored optical-flow-based residual propagation (utilizing both traditional Farneback and deep-learning RAFT models). Our rigorous failure analysis across the \verb|wildvideo| and \verb|bmx-trees| datasets reveals the fundamental constraints of using 2D optical flow to enforce 3D temporal consistency, concluding that joint-temporal diffusion generation is the necessary paradigm for future iterations.
\end{abstract}

\section{Introduction}
Current state-of-the-art video inpainting frameworks like ProPainter excel at flow-guided propagation and spatio-temporal matching. Nevertheless, when confronted with large occlusions or structurally complex backgrounds, they are fundamentally limited by the information available within the source video, often yielding blurry or repetitive patches.

To overcome this, we designed a hybrid pipeline wherein SDXL is deployed exclusively on structurally difficult sub-regions (extracted via a \verb|diffusion_target| mask) across sparsely sampled keyframes. While the generated textures exhibit exceptional individual quality, directly grafting them back into the ProPainter-restored video induces severe spatial tearing and temporal flickering. This report outlines our structured engineering approach to diagnosing and resolving these integration artifacts.

\section{Methodology}

\subsection{Automated Object Segmentation utilizing SAM2}
% TODO: Automation of prompt generation for SAM2 across video sequences.
% [This section is left blank for future implementation]
\textit{This section is reserved for the automated bounding box and point prompt generation pipeline relying on Segment Anything Model 2 (SAM2), designed to track complex dynamic objects without human-in-the-loop intervention.}

\subsection{Resolving Spatial Artifacts: Distance Transform Alpha Blending}
\textbf{The Halo Artifact:} The initial implementation blended the SDXL-generated patches utilizing standard morphological operations (erosion/dilation paired with Gaussian blurring). This resulted in a near-binary hard mask. Because the SDXL VAE decoder introduces 1-2 pixels of high-frequency padding artifacts along latent boundaries during decoding, this hard-cut approach pasted a distinct white "halo" directly onto the original frame.

\textbf{Geometric Solution:} We discarded the morphological masking in favor of mathematically continuous Distance Transform-based Alpha Blending. For a binary mask $M$, the Euclidean distance transform $D(x, y)$ calculates the precise geometric distance from a pixel to the mask's boundary. The localized blending weight $\alpha(x, y)$ is formulated as a linear gradient:
\begin{equation}
    \alpha(x, y) = \text{clip}\left(\frac{D(x, y)}{R}, 0, 1\right)
\end{equation}
where $R$ is the predetermined feather radius. At the outer edge, $\alpha=0.0$ wholly preserves the pristine ProPainter background. Moving inwards, the weight linearly scales to $1.0$, completely obscuring the VAE border and achieving a pixel-perfect, seamless transition.

\subsection{Addressing Temporal Flickering: Optical Flow Residual Propagation}
\textbf{The Flickering Artifact:} Because SDXL generates structural textures independently per keyframe, regions oscillate between high-resolution diffusion patches and lower-resolution ProPainter backgrounds, yielding a severe temporal instability penalty.

\textbf{Tracking Residuals:} To suppress this flickering, we attempted to propagate the "high-definition texture residuals" (the difference between the SDXL frame and the ProPainter baseline) from keyframes to neighboring non-keyframes. This was achieved by warping the residuals along estimated motion trajectories. We evaluated two optical flow estimators:
\begin{itemize}
    \item \textbf{OpenCV Farneback:} A traditional dense optical flow algorithm.
    \item \textbf{RAFT:} A state-of-the-art deep neural network for optical flow computation.
\end{itemize}
To defend against chaotic warping near occlusions, the algorithm enforces a strict \textit{Overlap Ratio}. Warped patches failing this structural coherency check are automatically discarded.

\section{Experiments and Analysis}

\subsubsection{Sparse Diffusion Targeting}

A key observation in our early experiments is that directly applying diffusion enhancement to the entire hard mask leads to overuse of generative modification, which damages surrounding background consistency and introduces seam artifacts.
Therefore, we explicitly restrict diffusion to a sparse subset of difficult pixels instead of the whole candidate region.

Our pipeline first applies ProPainter as the main restoration module and then constructs a \texttt{diffusion\_target} mask for local enhancement.
The target mask is obtained from candidate regions after excluding borrowable areas, followed by confidence gating, morphological filtering, and connected-component cleanup.
When the target region becomes too large, we enforce an upper bound on the ratio between diffusion-target pixels and hard-mask pixels, so that only the most difficult pixels are retained for diffusion refinement.
The difficulty score combines low support fraction and high color inconsistency, thereby prioritizing pixels that are less reliably restored by propagation-based methods.

This sparse design is supported by quantitative evidence.
Here, \texttt{target/hard} denotes the ratio between the number of pixels selected for diffusion refinement and the number of pixels in the hard mask.
We report both the mean ratio over all evaluated samples and the 90th percentile (p90), which reflects the upper-end behavior in high-coverage cases.
Before introducing the ratio constraint, the average \texttt{target/hard} ratio was approximately \(0.915\), indicating that diffusion was effectively applied to almost the entire hard region.
After the sparse strategy was introduced, this ratio was reduced to \(0.35\) or \(0.25\), confirming that diffusion was restricted to a much smaller and more selective subset of pixels.

\begin{table}[h]
\centering
\caption{Effect of sparse targeting on diffusion coverage. Lower values indicate more selective diffusion refinement.}
\label{tab:sparse_target}
\begin{tabular}{lcc}
\toprule
Setting & target/hard mean & target/hard p90 \\
\midrule
Before sparsification & 0.915 & 0.994 \\
Sparse target (\(0.35\)) & 0.350 & 0.350 \\
Sparse target (\(0.25\)) & 0.250 & 0.250 \\
\bottomrule
\end{tabular}
\end{table}

Among these settings, the sparse-\(0.35\) configuration was considered the most balanced in terms of quality and stability, and was therefore kept as the default setting inside the diffusion-enhanced branch.

\begin{figure}[h]
    \centering
    \includegraphics[width=\linewidth]{results/part3/wildvideo/blending_comparison_labeled_horizontal.png}
    \caption{Sparse diffusion targeting. From left to right: hard mask, diffusion target, and enhanced result.}
    \label{fig:sparse_target}
\end{figure}

\subsubsection{Improved Boundary Blending}

Besides overuse of diffusion, another major artifact in our early results is the visible halo or bright boundary around refined regions.
Our analysis suggests that this issue comes from the interaction between VAE decoding artifacts and an overly rigid alpha blending scheme, which effectively pastes boundary noise back into the restored frame.

To alleviate this problem, we adopted a distance-transform-based alpha blending strategy.
Instead of using a hard or poorly shaped transition band, the new design assigns low alpha values near the outer boundary and gradually increases the replacement weight toward the interior of the refined region.
This creates a much smoother geometric transition between the enhanced patch and the original restored frame.

In practice, this modification substantially reduces halo artifacts and improves single-frame visual naturalness. It plummeted the outside change L1 mean from \(0.0091\) to \(0.0023\) (a 74\% reduction) and the seam change L1 mean from \(0.6508\) to \(0.2841\).
It is one of the clearest technical improvements obtained in Part 3.
Nevertheless, this improvement mainly benefits spatial quality on individual frames, and does not by itself resolve the temporal instability introduced by frame-wise diffusion enhancement.

\begin{figure}[h]
    \centering
    \includegraphics[width=\linewidth]{figures/part3_boundary_blending.png}
    \caption{Improved boundary blending. The distance-transform alpha reduces visible halo and seam artifacts.}
    \label{fig:boundary_blending}
\end{figure}

\subsubsection{Failure Analysis of Flow-based Temporal Propagation}

Although localized diffusion can improve difficult regions on key frames, it also introduces temporal inconsistency because enhanced frames may become noticeably sharper than their neighbors.
To reduce this flickering effect, we further explored flow-based residual propagation, where high-quality details generated on key frames are propagated to nearby frames using optical flow.

However, this direction did not become our final solution.
Both our experiments and the underlying task structure suggest that flow-based propagation is fundamentally weak for object removal and background reconstruction.
For instance, while using traditional OpenCV Farneback flow on the \texttt{wildvideo} dataset successfully propagated textures and increased the hard artifact gain to \(0.0695\), the temporal instability ratio remained high at \(3.724\). More critically, on the highly dynamic \texttt{bmx-trees} dataset, or when switching to the high-precision RAFT model, the propagation broke down completely (hard gain plummeting back to \(0.000\)).
The RAFT failure specifically highlighted severe engineering integration constraints, including forward/backward mapping mismatches and tensor normalization conflicts.

The fundamental reason is that this task is not merely about transporting visible pixels across time.
Instead, it often requires generating previously occluded background content that is never visible in adjacent frames.
Optical flow is effective for estimating correspondence between already visible appearances, but it is not designed to hallucinate missing content.

This mismatch becomes most severe near occlusion boundaries.
After object removal, the missing regions often involve disocclusion, one-to-many correspondence, and invalid backward mappings, making the inverse problem ill-posed.
Moreover, 2D flow models are easily corrupted by parallax, non-rigid motion, motion blur, and interpolation around high-frequency boundaries, which explains why seam and outside-region metrics are particularly sensitive to propagation strength.
In addition, propagation is recursive by nature, so even small alignment errors accumulate over time and manifest as temporal flicker.

Our quantitative results are consistent with this theoretical analysis.

\begin{table*}[t]
\centering
\small
\setlength{\tabcolsep}{4pt}
\caption{Effect of propagation matching and flow models on the gain--risk trade-off.}
\label{tab:flow_failure}
\begin{tabular}{lccccccc}
\toprule
Dataset & Method & Radius & Outside $\downarrow$ & Seam $\downarrow$ & Leakage $\downarrow$ & Hard Gain $\uparrow$ & Temporal Ratio $\downarrow$ \\
\midrule
\multirow{3}{*}{\texttt{wildvideo}} & No Propagation     & 0 & 0.002300 & 0.284100 & 0.008900 & 0.041900 & 3.720000 \\
                                    & OpenCV Farneback   & 5 & 0.002100 & 0.284000 & 0.008600 & \textbf{0.069500} & 3.724000 \\
                                    & RAFT Flow          & 5 & 0.003100 & 0.377000 & 0.012500 & 0.000000 & 3.907000 \\
\midrule
\multirow{2}{*}{\texttt{bmx-trees}} & OpenCV Farneback   & 5 & -        & 1.224000 & -        & 0.000000 & 0.998000 \\
                                    & RAFT Flow          & 5 & -        & 1.117000 & -        & 0.000000 & 0.984000 \\
\bottomrule
\end{tabular}
\end{table*}

Overall, our experiments indicate that the diffusion-enhanced branch can improve some difficult single-frame cases, especially after sparse targeting and improved alpha blending.
However, it does not consistently outperform the baseline in overall evaluation, mainly because temporal stability remains the dominant bottleneck.
Therefore, we keep the earlier pipeline as the default system and regard the diffusion branch as an exploratory extension with informative failure cases.

\section{Discussion and Future Direction}

Our exhaustive debugging cycle establishes a fundamental architectural ceiling: utilizing 2D optical flow as a post-processing bandage to stitch together independently generated 2D diffusion patches is structurally unsound for highly dynamic real-world videos. Even with perfectly aligned optical flow, 2D matrices simply cannot hallucinate or reconcile 3D parallax shifts and dynamic occlusions without introducing severe temporal distortion.

\section{Conclusion}
This report successfully resolves the spatial halo dilemma facing hybrid ProPainter-Diffusion frameworks through rigorous mathematical Distance Transform blending. Concurrently, by thoroughly chronicling the failure of multiple optical-flow residual propagation attempts, we permanently rule out 2D flow-tracking as a viable temporal stabilizer. Moving forward, the project must abandon post-hoc spatial stitching and pivot strictly to Joint Temporal Generation architectures (e.g., 3D-UNet, Temporal Attention native to the latent space) to guarantee intrinsic inter-frame physical consistency.

\end{document}
